% PCT-SOURCE: pdf=019 print=007 kind=prose id=p7-superselection-experiment
It should be emphasized that the above super-selection rules for \(Q,B\), and
\((-1)^F\), like all laws of physics, depend upon experiment. It is quite
unclear at the moment whether further super-selection rules exist. For
example, it may be that there are laws of lepton conservation which define
super-selection rules.

% PCT-SOURCE: pdf=019 print=007 kind=section-title id=sec-ch1-symmetry-operations
\subsection{Symmetry Operations}\label{sec:ch1-symmetry-operations}

% PCT-SOURCE: pdf=019 print=007 kind=prose id=p7-symmetry-definition
A \emph{symmetry operation} (sometimes called an \emph{invariance principle},
or simply a \emph{symmetry}) of a physical system is a correspondence which
yields for each physically realizable state \(\ket{\Phi}\), another,
\(\ket{\Phi'}\), such that all transition probabilities are preserved:

% PCT-SOURCE: pdf=019 print=007 kind=display id=eq-ch1-transition-probability
\begin{equation}
  \left|\braket{\Phi'}{\Psi'}\right|^2
  =\left|\braket{\Phi}{\Psi}\right|^2.
  \tag{1-1}\label{eq:ch1-transition-probability}
\end{equation}

% PCT-SOURCE: pdf=019 print=007 kind=prose id=p7-symmetry-mapping
It is assumed that the mapping \(\ket{\Phi}\to\ket{\Phi'}\) is one to one.
This means that as \(\ket{\Phi}\) runs over all physically realizable states,
so does \(\ket{\Phi'}\), and if \(\ket{\Phi}\) and \(\ket{\Psi}\) are
distinct, so are \(\ket{\Phi'}\) and \(\ket{\Psi'}\). An example of a symmetry
is the operator translating the system by the four-vector \(a\). This is
represented by an operator \(V(a)\) which is unitary [i.e.,
\(\braket{V\Phi}{V\Psi}=\braket{\Phi}{\Psi}\)]. Another example is the
\(\mathrm{CPT}\) operator \(\Theta\), which is anti-unitary [i.e.,
\(\braket{\Theta\Phi}{\Theta\Psi}=\overline{\braket{\Phi}{\Psi}}\)].
Incidentally, the operator \(\Theta\) interchanges the coherent sub-spaces
with opposite charge and baryon number. Clearly, both unitary and anti-unitary
operators satisfy \eqref{eq:ch1-transition-probability}. In fact, all mappings \(\ket{\Phi}\to\ket{\Phi'}\)
lead essentially to a unique transformation \(\ket{\Phi}\to\ket{\Phi'}\)
satisfying \eqref{eq:ch1-transition-probability}, and this transformation is either unitary or anti-unitary
(Ref.~1).

% PCT-SOURCE: pdf=019 print=007 kind=theorem id=thm-ch1-symmetry
\begin{theorem}[1-1]\label{thm:ch1-symmetry}
% PCT-SOURCE: pdf=019 print=007 kind=theorem-statement id=thm-ch1-symmetry-statement
Let \(\ket{\Phi}\to\ket{\Phi'}\) be a symmetry of a physical theory
satisfying the hypothesis of commutative super-selection rules.

% PCT-SOURCE: pdf=019 print=007 kind=theorem-statement id=thm-ch1-symmetry-invariant
If the symmetry leaves coherent subspaces invariant, then there exists in each
coherent subspace a unitary or anti-unitary operator \(V\) such that for all
physically realizable states of that subspace

% PCT-SOURCE: pdf=019 print=007 kind=display id=eq-ch1-symmetry-action
\begin{equation}
  \ket{\Phi'}=V\ket{\Phi}.
  \tag{1-2}\label{eq:ch1-symmetry-action}
\end{equation}

% PCT-SOURCE: pdf=019 print=007 kind=theorem-statement id=thm-ch1-symmetry-uniqueness
The operator \(V\) is uniquely determined up to a phase.

% PCT-SOURCE: pdf=019 print=007 kind=theorem-statement id=thm-ch1-symmetry-moved
If the symmetry does not leave coherent subspaces invariant, then restricted to
a coherent subspace it is a one-to-one mapping onto another coherent
subspace, unitary or anti-unitary and unique up to a phase.
\end{theorem}

% PCT-SOURCE: pdf=020 print=008 kind=prose id=p8-theorem-essence
We shall not prove the theorem here; its essence is the following. The
one-to-one ray correspondence \(\ket{\Phi}\to\ket{\Phi'}\) can be induced by
one of many different vector correspondences, \(\ket{\Phi}\to\ket{\Phi'}\) in
the underlying Hilbert space, but in general such a correspondence is neither
linear nor anti-linear; i.e., neither

% PCT-SOURCE: pdf=020 print=008 kind=display id=eq-ch1-linear-ray-map
\begin{equation*}
  \alpha\ket{\Phi'}+\beta\ket{\Psi'}
  =\bigl(\alpha\ket{\Phi}+\beta\ket{\Psi}\bigr)'
\end{equation*}

% PCT-SOURCE: pdf=020 print=008 kind=prose id=p8-antilinear-intro
nor

% PCT-SOURCE: pdf=020 print=008 kind=display id=eq-ch1-antilinear-ray-map
\begin{equation*}
  \overline{\alpha}\ket{\Phi'}+\overline{\beta}\ket{\Psi'}
  =\bigl(\alpha\ket{\Phi}+\beta\ket{\Psi}\bigr)'
\end{equation*}

% PCT-SOURCE: pdf=020 print=008 kind=prose id=p8-ray-map-conclusion
holds. What the theorem says is that when the vectors of a single coherent
subspace are considered there is a linear or anti-linear transformation unique
up to a phase (one or the other, not both!) which yields the given
correspondence between rays.

% PCT-SOURCE: pdf=020 print=008 kind=prose id=p8-final-symmetry-remark
A final remark about symmetry operations. Our definition clearly makes every
unitary and anti-unitary operator a symmetry operator and Theorem \ref{thm:ch1-symmetry} shows
that these are essentially the only such. How then can one system be more
symmetrical than another? The answer lies in the physical interpretation of
the operations. Consider for example a theory of two spinless particles with
coordinate and momentum operators \(q_1(t),q_2(t),p_1(t),p_2(t)\). Then the
mapping \(q_1(t)\to-q_1(t),\ p_1(t)\to-p_1(t)\) of the observables onto
themselves can always be defined, but only when the system possesses symmetry
under space inversion will the mapping be induced by a unitary operator \(V\)
according to

% PCT-SOURCE: pdf=020 print=008 kind=display id=eq-ch1-space-inversion
\begin{equation}
  Vq_j(t)V^{-1}=-q_j(t),\qquad
  Vp_j(t)V^{-1}=-p_j(t),\qquad j=1,2
  \tag{1-3}\label{eq:ch1-space-inversion}
\end{equation}

% PCT-SOURCE: pdf=020 print=008 kind=prose id=p8-space-inversion-time
with \(V\) independent of \(t\). On the other hand, it may be possible to
define a unitary operator, \(V\), which has the effect of a parity operator for
the center of mass motion and total momentum, even if the theory is not
invariant under space-inversion, i.e., even if there is no \(V\) satisfying
\eqref{eq:ch1-space-inversion}. For example, \(V\) might satisfy

% PCT-SOURCE: pdf=020 print=008 kind=display id=eq-ch1-exchange-parity
\begin{equation*}
\begin{aligned}
  Vq_1(t)V^{-1}&=-q_2(t), & Vq_2(t)V^{-1}&=-q_1(t),\\
  Vp_1(t)V^{-1}&=-p_2(t), & Vp_2(t)V^{-1}&=-p_1(t).
\end{aligned}
\end{equation*}

% PCT-SOURCE: pdf=020 print=008 kind=prose id=p8-then
Then

% PCT-SOURCE: pdf=020 print=008 kind=display id=eq-ch1-center-of-mass-parity
\begin{equation*}
\begin{aligned}
  V\bigl(q_1(t)+q_2(t)\bigr)V^{-1}
    &=-\bigl(q_1(t)+q_2(t)\bigr),\\
  V\bigl(p_1(t)+p_2(t)\bigr)V^{-1}
    &=-\bigl(p_1(t)+p_2(t)\bigr).
\end{aligned}
\end{equation*}

% PCT-SOURCE: pdf=020 print=008 kind=prose id=p8-boundary
This example makes clear that the statement that a physical system possesses
space-inversion symmetry only makes sense when one has specified what space
inversion is supposed to do to the observables of
